\scanpage{420}
\mspage{332}
% archive sheet number: 419

% upper left, beside lines 1–4
% slanted
\margin{\unsure{Maladroit} \unsure{on peut|il faut} traiter les deux cas b) et \unsure{c')|a')} en $=$ temps cf. bas de \unsure{pge|pg} 333}

% top, lines 1–2
Supposons d'abord que $\underline S$ satisfasse b)\textsuperscript{*}. Donc on sait que
% upper third, lines 3–4
\begingroup\csname tagsleft@true\endcsname
\[
\begin{aligned}
x\in\operatorname{Ob}\underline S
  &\Longleftrightarrow \deleted{\unsure{\underline S}}\ \underline S(\pi)_{/x}\ \text{est \unsure{cofinal|cofinale}}\\
  &\hphantom{\Longleftrightarrow}\text{dans }M_{\pi}^{\deleted{\unsure{x}}}/x
\end{aligned}
\tag*{(**)}
\]
\endgroup

% upper third, lines 5–6
Nous voulons déduire l'équivalence (*) de l'équivalence (**).

% upper–middle, lines 7–12
Prouvons $\Rightarrow$ dans (*), i.e. supposons $x\in\operatorname{Ob}\underline S$, et prouvons que $x$ est dans l'image essentielle de $\operatorname{Ind}_{\pi}\underline S(\pi)$. En vertu de (**), $\underline S(\pi)_{/x}$ est \inserted{une sous-catégorie (str. pleine)} cofinale dans $M_{\pi}/x$, donc comme $M_{\pi}/x$ est \unsure{grd|grde} devant $\pi$, $\underline S(\pi)_{/x}$ l'est aussi. \deleted{\unsure{Donc enfin|Donc ceci}} Comme on sait que
% middle, lines 13–14
\[
x=\varinjlim_{i\in M_{\pi}/x}x_i
\]
% middle–lower, lines 15–16
($\pi$ étant cardinal utile pour $M$), il s'ensuit \deleted{que} par cofinalité qu'on a aussi
% lower third, lines 17–18
\[
x=\varinjlim_{i\in\underline S(\pi)_{/x}}x_{\unsure{.|i}}, 
\]
% lower third, lines 19–21
et comme $\underline S(\pi)_{/x}$ est grd devant $\pi$, cela implique que $x$ appartient à l'image essentielle de $\operatorname{Ind}_{\pi}(\underline S(\pi))$.

% foot, last two lines
Prouvons $\Leftarrow$ dans (*), i.e. supposons que $x$ soit dans l'image essentielle de $\operatorname{Ind}_{\pi}(\underline S(\pi))$,
